\section{Results and Design Validation}

The results section is structured as design validation: each experiment tests a specific claim about the framework's architecture. We first establish that TPRS substantially reduces the paralog gap (headline), then validate each design decision through component ablation, and finally characterize what the learned router captures.

\subsection{Headline: Closing the Paralog Gap}

Table~\ref{tab:benchmark} (TPRS row) reports end-to-end alignment performance. On paralogs, TPRS achieves primary exon F1 of 0.79 (mouse), 0.77 (dog), 0.76 (chicken), and 0.57 (zebrafish), representing improvements of $+$0.36--0.46 over the best single alignment tool per species. The best-overlap F1 reaches 0.84 (mouse), 0.85 (dog), 0.80 (chicken), and 0.63 (zebrafish), recovering 71--90\% of the one-to-one performance ceiling: 87--90\% in mammals and 71\% in zebrafish. The zebrafish ceiling is lower primarily due to the greater evolutionary distance (mean human--zebrafish protein identity of $\sim$55\% vs $\sim$80\% for human--mouse), which limits C-stage alignment quality even with species-matched proteins.

The gains are concentrated in the paralog regime, where existing single-tool approaches are most vulnerable. On one-to-one orthologs, existing tools already achieve $>$0.87 F1, leaving negligible headroom. On paralogs, TPRS targets both diagnosed failure modes, producing gains that are largest in the species with the deepest paralog gap (zebrafish: best single tool 0.17 $\rightarrow$ TPRS 0.57, $+$0.40). This concentration of gain on the hardest cases confirms that TPRS addresses the diagnosed root causes rather than an orthogonal source of error—its improvement is a direct consequence of breaking through the annotation-free ceiling identified in Section~2.

\subsection{Design Validation via Component Ablation}

To validate that each architectural decision addresses its intended failure mode, we ablate TPRS by removing one component at a time (E2--E6). Table~\ref{tab:ablation} and Fig.~\ref{fig:ablation} present the results. We structure the interpretation around the design claims established in Section~3.1.

\begin{table}[t]
\centering
\caption{Pipeline component ablation (Phase~5, 4 species).}
Mean F1 with 95\% paired-bootstrap CI. $n = 445$-$489$ queries per species (80/20 group-shuffle split by query). \textbf{Bold} = mean $|\Delta| > 0.20$ (critical component).
\label{tab:ablation}
\scriptsize
\setlength{\tabcolsep}{3.5pt}
\begin{tabular}{lccccc}
\toprule
\multicolumn{6}{c}{\textbf{Panel A: best\_overlap F1}} \\
Experiment & Mouse & Dog & Chicken & Zebrafish & Mean $\Delta$ \\
\midrule
E1: Full TPRS & 0.836 (0.761, 0.903) & 0.854 (0.794, 0.909) & 0.802 (0.727, 0.872) & 0.635 (0.543, 0.727) & -- \\
E6: - last (A-stage) & 0.546 (0.457, 0.634) & 0.553 (0.471, 0.634) & 0.391 (0.307, 0.476) & 0.312 (0.223, 0.402) & \textbf{-0.332} \\
E6: - mmseqs (A-stage) & 0.566 (0.479, 0.654) & 0.559 (0.478, 0.638) & 0.390 (0.306, 0.475) & 0.304 (0.220, 0.392) & \textbf{-0.327} \\
E3: - Router (heuristic) & 0.460 (0.371, 0.552) & 0.601 (0.521, 0.678) & 0.545 (0.457, 0.629) & 0.262 (0.186, 0.343) & \textbf{-0.315} \\
E6: - diamond (A-stage) & 0.595 (0.507, 0.681) & 0.565 (0.484, 0.644) & 0.391 (0.307, 0.476) & 0.327 (0.239, 0.417) & \textbf{-0.313} \\
E6: - miniprot (A-stage) & 0.595 (0.507, 0.681) & 0.568 (0.485, 0.646) & 0.419 (0.336, 0.505) & 0.316 (0.230, 0.404) & \textbf{-0.307} \\
E2: - Species protein & 0.616 (0.529, 0.698) & 0.706 (0.635, 0.772) & 0.546 (0.462, 0.631) & 0.402 (0.311, 0.489) & \textbf{-0.215} \\
E4: - Co-ortholog recovery & 0.640 (0.547, 0.733) & 0.762 (0.681, 0.836) & 0.594 (0.499, 0.687) & 0.559 (0.460, 0.656) & -0.143 \\
E5: - C-stage features\textsuperscript{\dag} & 0.807 (0.727, 0.880) & 0.852 (0.786, 0.910) & 0.805 (0.729, 0.877) & 0.637 (0.542, 0.730) & -0.007 \\
\midrule
\multicolumn{6}{c}{\textbf{Panel B: primary F1}} \\
Experiment & Mouse & Dog & Chicken & Zebrafish & Mean $\Delta$ \\
\midrule
E1: Full TPRS & 0.793 (0.712, 0.874) & 0.771 (0.691, 0.846) & 0.760 (0.675, 0.841) & 0.575 (0.475, 0.673) & -- \\
E3: - Router (heuristic) & 0.247 (0.165, 0.336) & 0.451 (0.365, 0.537) & 0.429 (0.340, 0.518) & 0.222 (0.147, 0.301) & \textbf{-0.388} \\
E6: - last (A-stage) & 0.338 (0.249, 0.431) & 0.445 (0.358, 0.529) & 0.385 (0.299, 0.472) & 0.195 (0.122, 0.275) & \textbf{-0.384} \\
E6: - miniprot (A-stage) & 0.354 (0.266, 0.446) & 0.452 (0.366, 0.536) & 0.385 (0.299, 0.471) & 0.205 (0.131, 0.285) & \textbf{-0.376} \\
E6: - diamond (A-stage) & 0.385 (0.293, 0.480) & 0.460 (0.372, 0.545) & 0.367 (0.283, 0.454) & 0.215 (0.139, 0.295) & \textbf{-0.368} \\
E6: - mmseqs (A-stage) & 0.390 (0.299, 0.486) & 0.462 (0.375, 0.548) & 0.377 (0.292, 0.464) & 0.216 (0.141, 0.296) & \textbf{-0.363} \\
E2: - Species protein & 0.355 (0.265, 0.449) & 0.594 (0.509, 0.676) & 0.482 (0.395, 0.571) & 0.278 (0.196, 0.363) & \textbf{-0.298} \\
E4: - Co-ortholog recovery & 0.631 (0.539, 0.723) & 0.658 (0.569, 0.743) & 0.532 (0.432, 0.630) & 0.502 (0.400, 0.604) & -0.144 \\
E5: - C-stage features\textsuperscript{\dag} & 0.784 (0.702, 0.865) & 0.740 (0.656, 0.819) & 0.727 (0.635, 0.813) & 0.566 (0.466, 0.663) & -0.021 \\
\bottomrule
\end{tabular}
\vspace{2pt}
\textsuperscript{\dag}C-stage features show small aggregate $\Delta$ (redundancy with species indicator) but primary F1 drops consistently in all 4 species (range $-0.010$ to $-0.033$). See Table~S5 for conditional analysis on cross-species subsets.\\
\textsuperscript{\ddag}Leave-one-A-tool-out robustness (Exp~8): the best fixed-hybrid H1 verdict is preserved on $23/24$ LOO trials across the 3 H1-passing paralog cells (mouse 7/8, dog 8/8, zebrafish 8/8). Ablation $\Delta$ values are stable within $\pm 0.02$ across K $\in \{10, 50, 200\}$ and $\tau \in \{0.90, 0.95, 0.99\}$.
\end{table}

\subsubsection{Claim 1: Candidate generation requires tool diversity (E6)}

\textbf{Experiment:} Remove each A-stage tool individually (leaving 3 of 4 in Phase~1).

\textbf{Result:} Removing any single tool reduces best-overlap F1 by $\Delta = -0.31$ to $-0.33$ (mean across 4 species). Primary F1 drops similarly ($\Delta = -0.36$ to $-0.38$), because fewer candidate regions mean the router has fewer options to choose from---and in some cases the correct locus is no longer in the candidate pool.

\textbf{Interpretation:} The four tools are complementary: no single tool dominates, and each removal causes a similar-magnitude drop. The consistent $\Delta$ magnitude across all four removals validates the choice to include diverse seeding and extension strategies rather than relying on a single best tool.

\subsubsection{Claim 2: Candidate generation requires query-type diversity (E2)}

\textbf{Experiment:} Remove the species-matched protein query---all candidates use the human query only.

\textbf{Result:} Primary F1 drops by $\Delta = -0.30$ (mean across 4 species). Best-overlap F1 drops by only $-0.22$.

\textbf{Interpretation:} The asymmetry between primary and best-overlap $\Delta$ is diagnostic. Species-matched protein predominantly helps the router select correctly, not merely inflate alignment scores: without it, the router cannot distinguish a correct-locus, low-identity human alignment from a wrong-locus, moderate-identity human alignment. The species protein at the correct locus produces a distinctively high-identity signal (c\_identity $>$ 0.95) that the router learns to prefer. This validates the query-identity dimension as the mechanism addressing RC1 within the joint selection framework.

\subsubsection{Claim 3: Selection must be learned (E3)}

\textbf{Experiment:} Replace the XGBoost router with the always-primary heuristic (select the locus overlapping the primary ortholog).

\textbf{Result:} Primary F1 drops by $\Delta = -0.39$, the largest single-component effect. Best-overlap F1 drops by $-0.32$.

\textbf{Interpretation:} No fixed rule can substitute for learned joint selection. The always-primary heuristic fails because, in the paralog regime, the correct locus is often \emph{not} the primary ortholog---it may be a different paralog copy with distinct genomic coordinates. The router learns to identify the correct locus from features beyond simple genomic position, including alignment quality and multi-tool evidence. This validates that RC2 requires a learned model.

\subsubsection{Claim 4: Prior knowledge fills blind spots (E4)}

\textbf{Experiment:} Remove Ensembl Compara co-ortholog recovery (Phase~2).

\textbf{Result:} Mean best-overlap F1 drops by $\Delta = -0.14$, primary F1 by $-0.14$.

\textbf{Interpretation:} The modest mean $\Delta$ masks a catastrophic effect on the affected subset. For the $\sim$27\% of queries where no A-stage tool hits the primary locus, removing Phase~2 drops best-overlap F1 from 0.64 to 0.21---the correct locus is no longer in the candidate pool at all, and the router has zero chance of selecting it. This validates the prior-knowledge dimension: tool diversity alone is not sufficient for all proteins, and curated orthology data provides essential coverage for blind spots.

\subsubsection{Claim 5: C-stage features provide orthogonal signal (E5)}

\textbf{Experiment:} Remove C-stage features from the router (retaining only A-stage, geometry, query, and species-indicator features).

\textbf{Result:} Primary F1 drops by $\Delta = -0.02$, best-overlap F1 by $-0.01$ (mean across 4 species). The small aggregate magnitude reflects redundancy with the species indicator, which already captures within-species vs.\ cross-species identity differences. However, the drop is consistent in direction across all 4 species.

\textbf{Interpretation:} C-stage features are free---computed from the miniprot output that the pipeline produces regardless---and provide consistent small gains. Their conditional contribution is larger on cross-species subsets ($+$0.02--0.04 F1, Table~S5), where C-stage quality metrics help distinguish genuine alignments from spurious ones.

\begin{figure}[t]
    \centering
    \includegraphics[width=\columnwidth]{figures/fig4_ablation_forest.pdf}
    \caption{Design validation via component ablation. $\Delta$ exon F1 from E1 (Full TPRS) for ablation experiments E2--E6, across 4 species. Diamonds = mean $\Delta$. Gray band = negligible effect ($|\Delta| < 0.05$). Each ablation tests a specific design claim from Section~3.1.}
    \label{fig:ablation}
\end{figure}

\subsection{Router vs.\ Heuristic Baselines}

To ensure the router's advantage reflects learned feature interactions rather than a weak baseline, we compare XGBoost against five non-learned heuristics: always-primary, max-C-stage-composite, max-identity, max-tool-consensus, and max-coverage (Table~S4). The router outperforms the strongest heuristic (max-identity) by $+$0.36 primary F1 (mean across 4 species) and $+$0.29 best-overlap F1. This gap demonstrates that learned routing captures conditional patterns---for instance, preferring a locus on the primary chromosome with moderate identity over a high-identity locus on a different chromosome---that no single heuristic or linear combination of heuristics can express.

\subsection{Directional Consistency Across Species}

Across all 8 ablation experiments (E2--E6, best-overlap and primary) and all 4 species, every $\Delta$ has the same sign (Table~\ref{tab:ablation}). The probability of 32/32 same-sign outcomes under a null of random sign assignment is $p = 2^{-31} < 5\times 10^{-10}$. This directional consistency indicates that the design decisions validated here generalize across species, supporting the modular, species-agnostic architecture of TPRS. Preliminary cross-species router experiments (Exp~7 D7, training on 3 species and testing on the held-out 4th) suggest that paralog gains transfer across species, though with measurable degradation on the most evolutionarily distant pairs; systematic cross-species router evaluation remains future work (see Section~5).

\subsection{What the Router Learns}

Feature importance analysis (Table~S3) reveals a biologically interpretable decision hierarchy that generalizes across species. The top three predictors are \texttt{c\_identity} (mean gain 0.19), \texttt{c\_coverage} (0.12), and \texttt{on\_primary\_chrom} (0.11). The presence of a geometric feature (\texttt{on\_primary\_chrom}) among the top predictors is notable: it encodes a genomic proximity prior---the correct paralog often resides on the same chromosome as the primary ortholog---that no heuristic baseline encodes but that the router discovers from data. The router has thus learned a sensible biological rule: prefer loci that (1) produce high-identity spliced alignments, (2) cover the full protein, and (3) are genomically proximal to the known ortholog, with these preferences applied conditionally rather than as fixed thresholds.
